\documentclass[12pt,a4paper]{article}

\usepackage{cmap}
\usepackage[T2A]{fontenc}
\usepackage[utf8]{inputenc}
\usepackage[russian]{babel}
\usepackage{amsmath,amssymb}
\usepackage[a4paper,margin=2cm]{geometry}

\newcommand{\R}{\mathbb{R}}
\newcommand{\dd}{\,d}

\setlength{\parindent}{0pt}
\setlength{\parskip}{0.45em}
\sloppy

\begin{document}

\begin{center}
{\Large\bfseries Билет 8}
\end{center}

\noindent\fbox{%
  \begin{minipage}{\dimexpr\linewidth-2\fboxsep-2\fboxrule\relax}
  \normalfont
Геометрические приложения определенного интеграла: площадь криволинейной
трапеции, объем тела вращения, длина кривой. Площадь поверхности вращения (без
доказательства).
  \end{minipage}%
}

\section*{Краткий план}

Геометрические приложения определенного интеграла строятся через предельный
переход от простых фигур. Для криволинейной трапеции ее приближают снизу и
сверху клеточными множествами, меры которых равны нижней и верхней суммам
Дарбу. Для тела вращения аналогично используют цилиндры, поэтому объем равен
интегралу площадей поперечных сечений. Длина кривой определяется как супремум
длин вписанных ломаных; для \(C^1\)-параметризации она равна
\(\int |r'(t)|\,dt\). Площадь поверхности вращения определяется как предел
площадей поверхностей, полученных вращением вписанных ломаных; формула
приводится без доказательства.

\section*{Определения}

\textbf{Криволинейная трапеция.}
Пусть \(f\colon [a,b]\to\R\) неотрицательна. Множество
\[
  E=\{(x,y)\in\R^2:\ a\le x\le b,\ 0\le y\le f(x)\}
\]
называется криволинейной трапецией под графиком \(f\).

\textbf{Тело вращения.}
Пусть \(f\colon [a,b]\to\R\) неотрицательна. Множество
\[
  G=
  \{(x,y,z)\in\R^3:\ x\in[a,b],\ \sqrt{y^2+z^2}\le f(x)\}
\]
называется телом, полученным вращением графика \(f\) вокруг оси \(Ox\).

\textbf{Длина кривой.}
Пусть кривая задана непрерывной вектор-функцией
\[
  \Gamma=\{r(t):t\in[a,b]\}\subset\R^n.
\]
Для разбиения \(T=\{t_0,\dots,t_I\}\) отрезка \([a,b]\) вписанная ломаная
\(P_T\) имеет длину
\[
  |P_T|=\sum_{i=1}^I |r(t_i)-r(t_{i-1})|.
\]
Длиной кривой называется
\[
  |\Gamma|=\sup_T |P_T|.
\]
Если эта величина конечна, кривая называется спрямляемой.

\textbf{Поверхность вращения и ее площадь.}
Пусть \(f\colon [a,b]\to\R\) неотрицательна. Поверхность вращения графика
функции \(f\) вокруг оси \(Ox\) есть
\[
  Q=
  \{(x,y,z)\in\R^3:\ x\in[a,b],\ \sqrt{y^2+z^2}=f(x)\}.
\]
Для разбиения \(T=\{x_0,\dots,x_I\}\) вращаем вокруг оси \(Ox\) ломаную,
вписанную в график \(f\). Полученная поверхность \(Q_T\) состоит из боковых
поверхностей усеченных конусов, и ее площадь равна
\[
  S(Q_T)=
  \pi\sum_{i=1}^I
  \sqrt{(x_i-x_{i-1})^2+\bigl(f(x_i)-f(x_{i-1})\bigr)^2}\,
  \bigl(f(x_{i-1})+f(x_i)\bigr).
\]
Число \(S\) называется площадью поверхности вращения \(Q\), если
\[
  S=\lim_{\ell(T)\to0} S(Q_T).
\]

\section*{Теоремы}

\textbf{1. Площадь криволинейной трапеции.}
Если \(f\colon [a,b]\to\R\) интегрируема по Риману и неотрицательна на
\([a,b]\), то криволинейная трапеция
\[
  E=\{(x,y):a\le x\le b,\ 0\le y\le f(x)\}
\]
измерима по Жордану и
\[
  \mu(E)=\int_a^b f(x)\dd x.
\]

\textbf{2. Объем тела вращения.}
Если \(f\colon [a,b]\to\R\) интегрируема по Риману и неотрицательна на
\([a,b]\), то тело вращения
\[
  G=\{(x,y,z):x\in[a,b],\ \sqrt{y^2+z^2}\le f(x)\}
\]
измеримо по Жордану и
\[
  \mu(G)=\pi\int_a^b f^2(x)\dd x.
\]

\textbf{3. Длина гладкой кривой.}
Если кривая
\[
  \Gamma=\{r(t):t\in[a,b]\}\subset\R^n
\]
задана непрерывно дифференцируемой вектор-функцией \(r(t)\), то
\[
  |\Gamma|=\int_a^b |r'(t)|\dd t.
\]
В частности, для графика \(y=f(x)\), где \(f\in C^1[a,b]\),
\[
  |\Gamma|=\int_a^b \sqrt{1+\bigl(f'(x)\bigr)^2}\dd x.
\]

\textbf{4. Площадь поверхности вращения.}
Если \(f\in C^1[a,b]\) и \(f(x)\ge0\) на \([a,b]\), то площадь поверхности,
полученной вращением графика \(y=f(x)\) вокруг оси \(Ox\), существует и равна
\[
  S=2\pi\int_a^b f(x)\sqrt{1+\bigl(f'(x)\bigr)^2}\dd x.
\]
Эта теорема в билете требуется без доказательства.

\section*{Доказательства}

\textbf{Площадь криволинейной трапеции.}
Возьмем произвольное разбиение
\[
  T=\{a=x_0<x_1<\dots<x_I=b\}
\]
и обозначим
\[
  m_i=\inf_{x\in[x_{i-1},x_i]} f(x),
  \qquad
  M_i=\sup_{x\in[x_{i-1},x_i]} f(x).
\]
Построим клеточные множества
\[
  A_T=\bigcup_{i=1}^I
  \{(x,y):x_{i-1}\le x\le x_i,\ 0\le y\le m_i\},
\]
\[
  B_T=\bigcup_{i=1}^I
  \{(x,y):x_{i-1}\le x\le x_i,\ 0\le y\le M_i\}.
\]
Тогда
\[
  A_T\subset E\subset B_T.
\]
Их меры равны нижней и верхней суммам Дарбу:
\[
  \mu(A_T)=\sum_{i=1}^I m_i(x_i-x_{i-1})=s(f;T),
\]
\[
  \mu(B_T)=\sum_{i=1}^I M_i(x_i-x_{i-1})=S(f;T).
\]
Отсюда для нижней и верхней мер Жордана множества \(E\)
\[
  s(f;T)\le \underline\mu(E)\le \overline\mu(E)\le S(f;T).
\]
Так как \(f\) интегрируема, при \(\ell(T)\to0\)
\[
  s(f;T)\to J,\qquad S(f;T)\to J,
  \qquad J=\int_a^b f(x)\dd x.
\]
Переходя к пределу, получаем
\[
  J\le \underline\mu(E)\le \overline\mu(E)\le J.
\]
Значит \(\underline\mu(E)=\overline\mu(E)=J\), то есть \(E\) измерима и
\(\mu(E)=\int_a^b f(x)\dd x\).

\textbf{Объем тела вращения.}
Для того же разбиения \(T\) и тех же \(m_i,M_i\) рассмотрим цилиндры
\[
  q_i=\{(x,y,z):x\in[x_{i-1},x_i],\ \sqrt{y^2+z^2}\le m_i\},
\]
\[
  Q_i=\{(x,y,z):x\in[x_{i-1},x_i],\ \sqrt{y^2+z^2}\le M_i\}.
\]
Положим \(A_T=\bigcup q_i\), \(B_T=\bigcup Q_i\). Тогда
\[
  A_T\subset G\subset B_T.
\]
Так как площадь круга радиуса \(r\) равна \(\pi r^2\), то
\[
  \mu(q_i)=\pi m_i^2(x_i-x_{i-1}),
  \qquad
  \mu(Q_i)=\pi M_i^2(x_i-x_{i-1}).
\]
Обозначим \(\varphi(x)=\pi f^2(x)\). Поскольку \(f\) интегрируема по Риману,
функция \(\varphi\) также интегрируема. Для неотрицательной \(f\)
\[
  \mu(A_T)=s(\varphi;T),\qquad \mu(B_T)=S(\varphi;T),
\]
и поэтому
\[
  s(\varphi;T)\le
  \underline\mu(G)\le \overline\mu(G)\le S(\varphi;T).
\]
При \(\ell(T)\to0\) обе суммы Дарбу стремятся к
\[
  \int_a^b \varphi(x)\dd x=\pi\int_a^b f^2(x)\dd x.
\]
Следовательно, \(G\) измеримо и
\[
  \mu(G)=\pi\int_a^b f^2(x)\dd x.
\]

\textbf{Длина кривой.}
Сначала напомним вспомогательный факт. Если \(r\in C^1[a,b]\), то кривая
спрямляема, а переменная длина дуги
\[
  s(t)=|\{r(u):u\in[a,t]\}|
\]
дифференцируема и
\[
  s'(t)=|r'(t)|.
\]
Действительно, для малого \(\Delta t>0\) обозначим
\[
  \Delta s=s(t+\Delta t)-s(t),\qquad
  \Delta r=r(t+\Delta t)-r(t).
\]
Длина хорды не превосходит длины дуги, поэтому
\[
  |\Delta r|\le \Delta s.
\]
С другой стороны, достаточное условие спрямляемости для \(C^1\)-кривой дает
\[
  \Delta s
  \le
  \max_{\xi\in[t,t+\Delta t]} |r'(\xi)|\,\Delta t.
\]
После деления на \(\Delta t\):
\[
  \left|\frac{\Delta r}{\Delta t}\right|
  \le
  \frac{\Delta s}{\Delta t}
  \le
  \max_{\xi\in[t,t+\Delta t]} |r'(\xi)|.
\]
При \(\Delta t\to+0\) левая и правая части стремятся к \(|r'(t)|\), значит
по теореме о трех функциях \(s'_+(t)=|r'(t)|\). Левые производные
получаются так же.

Теперь по формуле Ньютона--Лейбница
\[
  |\Gamma|=s(b)-s(a)=\int_a^b s'(t)\dd t
  =
  \int_a^b |r'(t)|\dd t.
\]
Если \(\Gamma\) является графиком \(y=f(x)\), то
\[
  r(x)=(x,f(x)),\qquad |r'(x)|=\sqrt{1+\bigl(f'(x)\bigr)^2},
\]
и поэтому
\[
  |\Gamma|=\int_a^b \sqrt{1+\bigl(f'(x)\bigr)^2}\dd x.
\]

\textbf{Площадь поверхности вращения.}
По условию билета доказательство этой формулы не требуется. Важно помнить
определение: площадь поверхности вращения есть предел площадей поверхностей
усеченных конусов, полученных вращением вписанных ломаных. При \(f\in C^1\)
этот предел существует и равен
\[
  2\pi\int_a^b f(x)\sqrt{1+\bigl(f'(x)\bigr)^2}\dd x.
\]

\section*{Финальная устная версия}

Для неотрицательной интегрируемой функции \(f\) криволинейная трапеция
\[
  E=\{(x,y):a\le x\le b,\ 0\le y\le f(x)\}
\]
измерима, и ее площадь равна
\[
  \mu(E)=\int_a^b f(x)\dd x.
\]
Доказательство такое: по разбиению \(T\) строим клеточные приближения
\(A_T\subset E\subset B_T\). Их меры равны нижней и верхней суммам Дарбу
\(s(f;T)\) и \(S(f;T)\). При измельчении разбиения обе суммы стремятся к
интегралу, значит нижняя и верхняя меры Жордана \(E\) совпадают с этим
интегралом.

Тело вращения вокруг оси \(Ox\) задается условием
\[
  G=\{(x,y,z):x\in[a,b],\ \sqrt{y^2+z^2}\le f(x)\}.
\]
Его сечение плоскостью \(x=\mathrm{const}\) есть круг радиуса \(f(x)\), поэтому
ожидаемая площадь сечения равна \(\pi f^2(x)\). Строгое доказательство
повторяет доказательство для трапеции: приближаем \(G\) цилиндрами радиусов
\(m_i\) и \(M_i\), получаем суммы Дарбу для \(\varphi(x)=\pi f^2(x)\). Поэтому
\[
  \mu(G)=\pi\int_a^b f^2(x)\dd x.
\]

Длина кривой
\[
  \Gamma=\{r(t):t\in[a,b]\}
\]
определяется как супремум длин вписанных ломаных:
\[
  |\Gamma|=\sup_T\sum_i |r(t_i)-r(t_{i-1})|.
\]
Если \(r\in C^1[a,b]\), то переменная длина дуги \(s(t)\) удовлетворяет
\(s'(t)=|r'(t)|\). Это следует из оценок
\[
  \left|\frac{\Delta r}{\Delta t}\right|
  \le
  \frac{\Delta s}{\Delta t}
  \le
  \max |r'|.
\]
По формуле Ньютона--Лейбница
\[
  |\Gamma|=\int_a^b |r'(t)|\dd t.
\]
Для графика \(y=f(x)\) это дает
\[
  |\Gamma|=\int_a^b\sqrt{1+\bigl(f'(x)\bigr)^2}\dd x.
\]

Наконец, площадь поверхности вращения определяется как предел площадей
поверхностей, полученных вращением вписанных ломаных. Если
\(f\in C^1[a,b]\), \(f\ge0\), то без доказательства используем формулу
\[
  S=2\pi\int_a^b f(x)\sqrt{1+\bigl(f'(x)\bigr)^2}\dd x.
\]

\section*{Источники}

Иванов Г.Е. \emph{Лекции по математическому анализу. Часть 1}:
\texttt{IvGE\_1.pdf}, стр. 230--238, 146--150.

Основные роли источника: стр. 230--231 -- площадь криволинейной трапеции;
стр. 232--233 -- объем тела вращения; стр. 234--237 -- определение и формула
площади поверхности вращения; стр. 237--238 -- формула длины дуги как
геометрическое приложение интеграла; стр. 146--150 -- определение длины кривой,
переменная длина дуги и доказательство \(s'(t)=|r'(t)|\).

Дополнительно был просмотрен \texttt{IvGE\_2.pdf}: прямого материала для
этого билета там нет; найденные страницы 171--173 относятся к более позднему
общему изложению риманова объема, длины и площади на многообразиях.

\end{document}
